\documentclass[11pt,letterpaper]{article}
\usepackage{physical_ai_legislation}
\usepackage[backend=biber,style=numeric,sorting=none]{biblatex}
\addbibresource{physical_ai_parking.bib}

\title{%
  \textbf{Physical AI Oncology Clinical Trial Site}\\
  \textbf{Parking and Patient Transportation Standards}\\[6pt]
  \large Parking Facility Design, Patient Access,\\
  and Transportation Demand Management\\[6pt]
  \normalsize Standards for the dedicated parking facility and patient\\
  transportation infrastructure serving Physical AI trial sites.
}
\author{%
  Kevin Kawchak\\
  CEO, ChemicalQDevice\\
  \texttt{kevink@chemicalqdevice.com}
}
\date{March 24, 2026}

\begin{document}
\maketitle
\thispagestyle{fancy}

\begin{center}
\footnotesize
\textit{This draft is independent and is not endorsed, sponsored, or approved
by any trial sponsor, CRO, site, IRB, regulator, or medical society;
and was adapted using Claude Code Opus 4.6.}
\end{center}

\vspace{1em}
\tableofcontents
\newpage

\section{Scope and Purpose}

These standards govern the design, construction, and operation of dedicated
parking facilities and patient transportation infrastructure for Physical AI
oncology clinical trial sites authorized in California. The parking facility
is a new-construction structure purpose-built to serve the trial site, designed
to accommodate the unique patient flow patterns of a 24-hour, on-demand
oncology clinical trial.

The 24-hour simulation demonstrated arrival patterns ranging from 2 patients
during overnight hours (hours 02 through 04) to 15 patients during peak hour
(hour 09), with a total of 168 unique patients in 24 hours and a peak
concurrent on-site population of 28 patients~\cite{newtrial2026}. The parking
facility must accommodate these variable demand patterns while maintaining the
8-minute average wait time and the 24/7 accessibility that define the Physical
AI trial format~\cite{formatcomparison2026}.

\section{Parking Facility Design}

\subsection{Capacity Requirements}

\begin{enumerate}[label=(\alph*)]
\item The parking facility shall provide a minimum of 120 parking spaces,
  calculated as follows:
  \begin{enumerate}[label=(\arabic*)]
  \item 80 patient and visitor spaces, based on the peak concurrent on-site
    population of 28 patients plus accompanying persons and a turnover factor
    reflecting variable procedure durations (15 minutes for imaging to 180
    minutes for surgery)~\cite{newtrial2026,patientinstructions2026}.
  \item 15 staff spaces for the 8 to 10 full-time equivalent staff plus
    on-call personnel across three shifts~\cite{sitespec2026}.
  \item 10 accessible spaces (8.3 percent of total), exceeding the minimum
    ADA and California Building Code Chapter 11B requirements for medical
    facilities.
  \item 5 short-term patient drop-off and pick-up spaces adjacent to the
    building entrance.
  \item 10 spaces designated for ride-share and medical transport vehicle
    staging.
  \end{enumerate}

\item The facility shall accommodate future expansion of up to 50 percent
  additional capacity through structural design that supports the addition
  of one parking level.
\end{enumerate}

\subsection{Structure and Configuration}

\begin{enumerate}[label=(\alph*)]
\item The parking facility shall be a new-construction structure, which may be
  either a surface lot or a multi-level structure (not to exceed three levels
  above grade), depending on the site constraints and zoning requirements.

\item If a multi-level structure is constructed, all patient parking levels
  shall be served by a minimum of two elevators and one ramp system meeting
  ADA accessibility standards.

\item The facility shall provide a covered pedestrian connection from the
  parking structure to the clinical building entrance, with a maximum walking
  distance of 300 feet from the most distant parking space to the building
  entrance.

\item Clear ceiling height shall be a minimum of 8 feet 2 inches on all
  levels to accommodate medical transport vehicles and passenger vans.

\item Drive aisles shall be a minimum of 24 feet wide for two-way traffic,
  with parking spaces a minimum of 9 feet wide by 18 feet deep.
\end{enumerate}

\subsection{Accessible Parking}

\begin{enumerate}[label=(\alph*)]
\item Accessible parking spaces shall be located at the level closest to the
  building entrance, with a maximum distance of 100 feet from the space to
  the building entry.

\item A minimum of two van-accessible spaces (8 feet wide with 8-foot access
  aisle) shall be provided.

\item Accessible routes from parking to the building entrance shall be covered,
  illuminated to a minimum of 10 foot-candles, and free of steps, steep grades
  (maximum slope 1:20), or other barriers.

\item Accessible parking signage shall include the international symbol of
  accessibility and shall note that the facility operates 24 hours per day
  for on-demand oncology care.
\end{enumerate}

\section{Patient Drop-Off Zone}

\begin{enumerate}[label=(\alph*)]
\item A dedicated patient drop-off and pick-up zone shall be located immediately
  adjacent to the main building entrance, separated from the parking facility
  access drive.

\item The drop-off zone shall accommodate a minimum of three vehicles
  simultaneously under a covered canopy extending at least 20 feet from the
  building facade.

\item The drop-off zone shall include a level, non-slip surface, curb cuts,
  and seating for patients waiting for transportation.

\item A call button or intercom shall allow patients to request assistance
  from site safety officers for transfer from vehicles to the building
  entrance.

\item The drop-off zone shall be operational 24 hours per day, with adequate
  lighting (minimum 10 foot-candles) and weather protection for all hours of
  operation, consistent with the simulation's demonstrated overnight patient
  arrivals and departures~\cite{newtrial2026}.
\end{enumerate}

\section{Transportation Demand Management}

\subsection{Public Transit Integration}

\begin{enumerate}[label=(\alph*)]
\item The site shall be located within one-quarter mile of a BART station or
  Muni stop, or the site operator shall provide a shuttle service connecting
  the site to the nearest transit hub at intervals not exceeding 30 minutes
  during all operating hours.

\item Real-time transit information displays shall be installed in the patient
  waiting area and the parking facility lobby, showing schedules for BART,
  Muni, and site shuttle services.

\item The site shall participate in the Clipper card program or equivalent
  regional transit payment system to facilitate patient transit use.
\end{enumerate}

\subsection{Ride-Share and Medical Transport}

\begin{enumerate}[label=(\alph*)]
\item The site shall designate a minimum of 10 parking spaces for ride-share
  and medical transport vehicle staging, located near the patient drop-off
  zone but not blocking drop-off lane traffic flow.

\item The site shall coordinate with ride-share services to establish a
  designated pick-up and drop-off geofence that directs drivers to the
  correct facility entrance.

\item For patients who require medical transport following procedures (such as
  those recovering from anesthesia after surgical robot procedures lasting
  90 to 180 minutes), the site shall coordinate discharge scheduling with
  transport providers to minimize patient wait
  times~\cite{patientinstructions2026}.
\end{enumerate}

\subsection{Bicycle and Pedestrian Access}

\begin{enumerate}[label=(\alph*)]
\item The site shall provide a minimum of 20 Class I (enclosed and secured)
  bicycle parking spaces and 10 Class II (outdoor rack) bicycle parking spaces.

\item Bicycle parking shall be located within 100 feet of the building entrance
  and shall include electric bicycle charging capability for at least 4 spaces.

\item Pedestrian access from the public sidewalk to the building entrance shall
  be clearly marked, well-lit (minimum 5 foot-candles), and separated from
  vehicle traffic.
\end{enumerate}

\section{24-Hour Parking Operations}

\subsection{Staffing and Monitoring}

\begin{enumerate}[label=(\alph*)]
\item The parking facility shall be monitored 24 hours per day through video
  surveillance and emergency call stations, consistent with the 24-hour
  operating schedule of the clinical facility.

\item Emergency call stations shall be located at each elevator lobby, at each
  stairwell entrance, and at intervals not exceeding 200 feet in the parking
  area.

\item A parking management system shall track occupancy in real time and display
  available spaces at the facility entrance, directing arriving patients to
  available spaces efficiently.

\item During overnight hours (10:00 PM to 6:00 AM), the parking facility may
  reduce lighting to 75 percent of daytime levels in unoccupied areas while
  maintaining full lighting in occupied areas, pedestrian routes, and the
  patient drop-off zone.
\end{enumerate}

\subsection{Electric Vehicle Support}

\begin{enumerate}[label=(\alph*)]
\item A minimum of 20 percent of parking spaces (24 spaces) shall be equipped
  with Level 2 electric vehicle charging stations (240V, 32A minimum).

\item A minimum of 4 parking spaces shall be equipped with Level 3 DC fast
  charging stations (minimum 50 kW), enabling patients to charge during shorter
  procedures such as imaging (10 to 20 minutes) or companion robot sessions
  (10 to 20 minutes)~\cite{patientinstructions2026}.

\item An additional 20 percent of parking spaces (24 spaces) shall be designed
  as EV-capable, with conduit and panel capacity installed during construction
  to support future charging station installation.

\item EV charging shall be available 24 hours per day, consistent with the
  facility's operating schedule.
\end{enumerate}

\section{Emergency Vehicle Access}

\begin{enumerate}[label=(\alph*)]
\item A dedicated ambulance bay shall be provided separate from both the parking
  facility and the patient drop-off zone, accommodating a minimum of two
  ambulances simultaneously.

\item Fire department access lanes shall be maintained along the full perimeter
  of both the clinical building and the parking facility, with a minimum width
  of 26 feet and a minimum vertical clearance of 13 feet 6 inches.

\item Emergency vehicle access routes shall be clearly marked and kept
  unobstructed at all times, with automatic enforcement through surveillance
  and notification systems.

\item The parking facility design shall not impede emergency vehicle access to
  the clinical building from any direction.
\end{enumerate}

\section{Sustainability}

\begin{enumerate}[label=(\alph*)]
\item The parking facility shall include a stormwater management system
  compliant with the San Francisco Green Infrastructure Ordinance,
  incorporating bioretention areas, permeable pavement, or equivalent
  green infrastructure techniques.

\item Rooftop areas of the parking structure, if applicable, shall incorporate
  either solar photovoltaic panels or a green roof system covering a minimum
  of 50 percent of the roof area.

\item Parking facility lighting shall use LED fixtures with occupancy sensors
  and daylight harvesting to minimize energy consumption during the 24-hour
  operating period.

\item Construction materials shall meet the CalGreen (Title 24, Part 11)
  requirements for recycled content and low-emission materials.
\end{enumerate}

\printbibliography[title={References}]

\end{document}
